\section{Parallel Reachability with a Shared Visited Set}\label{sec:approach}

Our starting point is the observation that a parallel timed checker spends its effort in two
places --- generating successors and inserting them --- and that only the first is naturally
parallel. How the second is organised determines everything else, so we describe the design as
a sequence of decisions, each forced by the previous one.

\subsection{One level of the search}

The search is level-synchronous. \Cref{fig:level} shows a single iteration: the frontier holds
the symbolic states discovered at level $k$, threads take slices of it under a dynamic
schedule, and each thread expands its slice and inserts the resulting successors into the
shared visited set. States that survive insertion are stamped with the level number and form
the frontier of level $k{+}1$; a barrier closes the level.

\begin{figure}[t]
\centering
\begin{tikzpicture}[font=\small,x=1mm,y=1mm]
\node[cpubox, minimum width=24mm, minimum height=13mm] (fr) at (0,0)
  {\textbf{frontier} $F_k$\\[1pt]\scriptsize states found at level $k$};
\node[cpubox, minimum width=30mm] (t1) at (46,11) {expand \textbf{and} insert};
\node[cpubox, minimum width=30mm] (t2) at (46,0)  {expand \textbf{and} insert};
\node[cpubox, minimum width=30mm] (t3) at (46,-11){expand \textbf{and} insert};
\node[lbl] at (46,19) {\ttfamily \#pragma omp for schedule(dynamic)};
\node[cpubox, fill=black!6, draw=black!55, minimum width=26mm, minimum height=30mm] (tab) at (92,0)
  {shared visited set\\[2pt]\scriptsize striped: one lock per stripe};
\node[cpubox, minimum width=24mm, minimum height=13mm] (nx) at (132,0)
  {\textbf{frontier} $F_{k+1}$\\[1pt]\scriptsize survivors of level $k$};
\foreach \t in {t1,t2,t3} { \draw[flow] (fr) -- (\t); \draw[flow] (\t) -- (tab); }
\draw[flow] (tab) -- node[lbl,above]{stamped $k$} (nx);
\draw[flow, dashed] (nx.south) -- ++(0,-20)
  -- node[lbl,above]{\textbf{barrier}: level $k$ ends} ($(fr.south)+(0,-20)$) -- (fr.south);
\end{tikzpicture}
\caption{One level of the search. Each thread expands \emph{and} inserts its own successors,
so no successor is ever handed between threads. Everything here is standard breadth-first
search except the box on the right, which is where the design question lies.}
\label{fig:level}
\end{figure}

Level synchronisation costs one barrier per level, of which there are typically a few hundred,
and buys two properties we rely on. Membership of the current level becomes a single integer
comparison against the stamp, which is how survivors are identified without an auxiliary set.
And there is a well-defined point at which per-level work can be scheduled, which
\cref{sec:batching} exploits.

\begin{figure*}[t]
\centering
\begin{tikzpicture}[font=\small, x=1mm, y=1mm,
  ppanel/.style={draw=black!45, rounded corners=2pt, fill=black!3},
  parr/.style={draw=black!55, rounded corners=1pt, fill=white, minimum height=6mm,
              font=\ttfamily\scriptsize, inner xsep=3pt},
  pstep/.style={draw=black!55, rounded corners=1pt, fill=black!8, minimum height=5.5mm,
               font=\scriptsize, inner xsep=3pt},
  pthr/.style={draw=black!60, rounded corners=1pt, fill=white, minimum width=8mm,
              minimum height=5.5mm, font=\ttfamily\scriptsize},
  pchain/.style={draw=black!55, fill=white, minimum width=7mm, minimum height=5mm,
                font=\scriptsize},
  pdead/.style={pchain, fill=black!12, text=black!45},
  pnote/.style={font=\scriptsize\itshape, text=black!60},
  pfl/.style={-{Latex[length=1.7mm]}, draw=black!65}]

% ================= LEFT: frontier =================
\node[ppanel, minimum width=44mm, minimum height=52mm, anchor=north west] (P1) at (0,0) {};
\node[font=\bfseries\small, anchor=north] at (22,-2) {FRONTIER $F_k$};
\node[pnote, anchor=north] at (22,-8) {flat, struct-of-arrays};
\node[parr, anchor=north west, minimum width=36mm] at (4,-14) {fk[\,] -- packed 64-bit keys};
\node[parr, anchor=north west, minimum width=36mm] at (4,-23) {fz[\,] -- zone indices (4 B)};
\node[pnote, anchor=north west, text width=36mm] at (4,-32)
  {the discrete configuration is \emph{recovered} from the key, never stored beside a zone};
\node[pnote, anchor=north west, text width=36mm] at (4,-44)
  {one slice per thread, \texttt{schedule(dynamic)}};

% ================= MIDDLE: one thread =================
\node[ppanel, minimum width=62mm, minimum height=52mm, anchor=north west] (P2) at (50,0) {};
\node[font=\bfseries\small, anchor=north] at (81,-2) {WHAT ONE THREAD DOES};
\foreach \i/\x in {0/54, 1/64, 2/74, 3/84, 4/94, 5/104} \node[pthr] at (\x,-10) {t\i};
\node[pnote, anchor=west] at (54,-16) {each takes a slice of $F_k$; no successor is ever handed over};

\node[pstep, anchor=north west, minimum width=52mm] at (56,-20)
  {\textbf{1} \texttt{gen\_succ\_reuse}: build successor \emph{in place}};
\node[pstep, anchor=north west, minimum width=52mm] at (56,-27)
  {\textbf{2} $\dk \gets$ pack(locations, integers) \quad 64-bit};
\node[pstep, anchor=north west, minimum width=52mm] at (56,-34)
  {\textbf{3} $h \gets \dk\cdot\varphi^{-1}$: \emph{high} bits $\to$ stripe, \emph{low} $\to$ slot};
\node[pstep, anchor=north west, minimum width=52mm, fill=black!18] at (56,-41)
  {\textbf{4} \textsc{lock}(stripe) $\to$ \textsc{Insert} $\to$ \textsc{unlock}};
\node[pnote, anchor=north west, text width=52mm] at (56,-48)
  {no private buffer, no routing pass: the producer stores it};

% ================= RIGHT: the table =================
\node[ppanel, minimum width=58mm, minimum height=52mm, anchor=north west] (P3) at (118,0) {};
\node[font=\bfseries\small, anchor=north] at (147,-2) {SHARED VISITED SET};
\node[pnote, anchor=north] at (147,-8) {$N_B \ge \max(1024,\,64p)$ stripes, one mutex each};

\foreach \s/\y in {0/-14, 1/-23, 2/-32} {
  \node[draw=black!55, fill=black!10, minimum width=7mm, minimum height=5mm,
        font=\scriptsize, anchor=west] (L\s) at (122,\y) {\s};
  \node[pnote, anchor=east] at (122,\y) {};
}
\node[pchain, anchor=west] (a1) at (133,-14) {$z_5$};
\node[pchain, anchor=west] (a2) at (141,-14) {$z_3$};
\node[pdead,  anchor=west] (a3) at (149,-14) {$z_1$};
\draw[pfl] (L0) -- (a1); \draw[pfl] (a1) -- (a2); \draw[pfl] (a2) -- (a3);
\node[pchain, anchor=west] (b1) at (133,-23) {$z_9$};
\draw[pfl] (L1) -- (b1);
\node[pchain, anchor=west] (c1) at (133,-32) {$z_7$};
\node[pchain, anchor=west] (c2) at (141,-32) {$z_4$};
\draw[pfl] (L2) -- (c1); \draw[pfl] (c1) -- (c2);
\node[pnote, anchor=north west, text width=52mm] at (122,-38)
  {an antichain per configuration, newest first; a subsumed zone is flagged \emph{dead}, never
   unlinked, so a thread reading one chain is never disturbed by another};

% ================= flows between panels =================
\draw[pfl, line width=0.5pt] (44,-26) -- (50,-26);
\draw[pfl, line width=0.5pt] (112,-26) -- (118,-26);

% ================= bottom: survivors + inset =================
\node[ppanel, minimum width=176mm, minimum height=26mm, anchor=north west] (P4) at (0,-56) {};
\node[font=\bfseries\small, anchor=north west] at (4,-58) {\textsc{Insert}, under the stripe lock};
\node[anchor=north west, text width=86mm, font=\scriptsize] at (4,-64)
  {(a) is the newcomer subsumed by a stored zone? $\Rightarrow$ discard\\
   (b) does it subsume stored zones? $\Rightarrow$ flag them dead\\
   (c) is any union with a stored zone convex? $\Rightarrow$ absorb, repeat\\
   (d) intern the zone, prepend the node, stamp it with level $k$};
\node[anchor=north west, text width=80mm, font=\scriptsize\itshape, text=black!60] at (94,-64)
  {(a)--(c) each traverse the antichain, so the critical section is $\Theta(\rho)$ --- this is
   the term \cref{sec:batching} removes. Interning in (d) is what makes a stored state cost a
   4-byte index: Viking $N{=}15$ has $2\,768\,894$ symbolic states drawn from $120$ distinct
   matrices. The stamp in (d) makes ``survivor of level $k$'' an integer comparison, so the
   next frontier is gathered without an auxiliary set.};
\end{tikzpicture}
\caption{Expansion and insertion, end to end. The frontier is held as flat arrays of packed
keys and zone indices; threads take slices of it under a dynamic schedule and each thread runs
the whole pipeline on its own successors, inserting them itself. One multiplicative hash yields
both the stripe (high bits) and the slot within it (low bits), so a thread locks only the
stripe it touches while the other $p-1$ proceed. Because there is no per-thread buffer and no
routing pass, every successor is written exactly once. The inset lists what one thread executes
inside the lock.}
\label{fig:pipeline}
\end{figure*}

\subsection{Why insertion is fused into expansion}

The natural way to parallelise while avoiding locks is to partition: hash each discrete
configuration to a shard, give each shard to one thread, and let that thread own its slice of
the visited set outright. We built such a backend and it has a structural cost that is easy to
overlook. Because a thread may not write into another thread's shard, expansion cannot insert;
it can only \emph{accumulate}. Successors are written to a private buffer, a routing pass then
reads every successor and dispatches it to the owning shard, and only then can merging begin.
\Cref{fig:pipeline} shows the resulting pipeline in full, and \cref{fig:dataflow} isolates the difference from partitioned ownership. Routing is serial bookkeeping proportional
to the number of successors --- tens of millions over a run --- and the private buffers hold a
second complete copy of the level at exactly the moment memory peaks.

\begin{figure}[t]
\centering
\begin{tikzpicture}[font=\small]
\node[lbl] at (0,1.55) {\textbf{(a) partitioned ownership}};
\node[cpubox,minimum width=15mm] (afr) at (0,0.7) {frontier};
\node[cpubox,minimum width=15mm] (aex) at (2.4,0.7) {expand};
\node[box,fill=black!4,draw=black!35,minimum width=17mm] (abuf) at (4.9,0.7)
   {private buffers\\[-1pt]\scriptsize\emph{2nd copy of level}};
\node[box,fill=black!12,draw=black!60,minimum width=15mm] (art) at (7.4,0.7)
   {route\\[-1pt]\scriptsize\emph{serial}};
\node[cpubox,fill=black!6,draw=black!55,minimum width=15mm] (ash) at (9.8,0.7) {merge\\shards};
\draw[flow] (afr)--(aex); \draw[flow] (aex)--(abuf); \draw[flow] (abuf)--(art);
\draw[flow] (art)--(ash);
\node[lbl] at (0,-0.85) {\textbf{(b) fused (this paper)}};
\node[cpubox,minimum width=15mm] (bfr) at (0,-1.7) {frontier};
\node[cpubox,minimum width=26mm] (bex) at (3.2,-1.7)
   {expand \textbf{and} insert\\[-1pt]\scriptsize one pass, no buffer};
\node[cpubox,fill=black!6,draw=black!55,minimum width=24mm] (btab) at (7.4,-1.7)
   {shared visited set\\[-1pt]\scriptsize stripe lock per key};
\draw[flow] (bfr)--(bex); \draw[flow] (bex)--(btab);
\draw[flow,dashed] (btab.south) -- ++(0,-0.62)
   -- node[lbl,below]{survivors $=$ next frontier} ($(bfr.south)+(0,-0.62)$) -- (bfr.south);
\end{tikzpicture}
\caption{Partitioned ownership needs three writes per successor --- into a buffer, out of it
during routing, and into a shard --- because a thread may not write outside its own shard.
Fusing expansion and insertion removes the buffer and the routing pass, and each successor is
written once.}
\label{fig:dataflow}
\end{figure}

We therefore share the visited set and pay for it with locking, but at a granularity fine
enough that the payment is small. The set is an array of $N_B$ independent open-addressed
sub-tables, each with its own mutex. A discrete key selects its stripe from the high bits of a
multiplicative hash and its slot within that stripe from the low bits of the same product, so
one multiplication serves both and the two choices are uncorrelated. We take
$N_B \ge \max(1024, 64p)$ for $p$ threads: with at least sixty-four stripes per thread,
collisions between threads are rare, at a cost of a few kilobytes of headers. Within a stripe,
each configuration owns a singly linked chain of nodes $\langle$zone id, level, next,
dead$\rangle$ drawn from an arena, so chains are addressed by integer indices that survive
table growth, and removal is logical: a subsumed zone is flagged dead rather than unlinked,
so a thread traversing one chain is never disturbed by activity on another.

\subsection{Making sharing affordable}

Fusing insertion into expansion only pays if insertion is cheap in memory as well as in time,
and two representation choices do most of that work.

The first concerns the discrete configuration. Rather than store it beside every zone, we pack
it losslessly into a single 64-bit key: a fixed bit-field per process location and per integer
variable, each width derived from the number of locations or the declared range of the
variable. Keys are then hashed and compared as machine words, and --- more usefully --- the
configuration is \emph{recovered} from the key by the inverse map when it is needed, so it
costs nothing per stored state. The packing is injective only when integer variables carry
declared ranges; an unranged variable claims sixteen bits, and a model whose fields exceed the
budget is rejected with a diagnostic rather than silently aliased, since aliasing two distinct
configurations would be unsound.

The second concerns zones, and rests on a measurement. Across an entire Viking $N{=}15$ run,
$2\,768\,894$ symbolic states contain only $120$ \emph{distinct} DBMs. Storing an explicit
matrix per symbolic state therefore wastes almost all of the space it occupies. We store each
distinct zone once in a shared interning table and refer to it by a four-byte index, with a
per-thread cache in front of the shared structure so repeated zones are resolved without
touching it. Together with key packing, this is what produces the memory results of
\cref{sec:eval}, and it is not a constant-factor saving: it converts the dominant term from
the number of symbolic states to the number of distinct zones.

A third choice concerns time rather than space. A generator returning a vector of successors
allocates a location vector, an integer vector and a DBM for each one, all discarded moments
later by insertion, which interns the zone by copy. Ours instead builds each successor
\emph{in place} in a reused slot of a per-thread buffer and returns a count; slots survive to
be overwritten by the next frontier state, so after a brief warm-up per-successor allocation
is zero. On Viking $N{=}17$ this alone took a run from roughly twelve seconds to $8.2$.

\subsection{The insertion algorithm}

\Cref{alg:insert} gives insertion in full. The parameters are the successor $(\ell,\Zone)$,
its packed key $\dk$, and the current level $k$; the stripe lock for $\dk$ is held throughout.

\begin{algorithm}[t]
\caption{$\textsc{Insert}(\ell, \Zone, \dk, k)$ --- executed under the stripe lock of $\dk$}
\label{alg:insert}
\begin{algorithmic}[1]
\State $A \gets \textsc{Chain}(\dk)$ \Comment{antichain currently stored for this configuration}
\State $(L,U) \gets \textsc{Bounds}(\ell)$ \label{ln:bounds}
\ForAll{live $v \in A$} \label{ln:sub1}
  \If{$\Zone \aLU[L,U] \Zone_v$} \State \Return $\bot$ \Comment{redundant: nothing to do}\label{ln:drop}
  \EndIf
\EndFor
\ForAll{live $v \in A$} \label{ln:sub2}
  \If{$\Zone_v \aLU[L,U] \Zone$} \State $v.\mathit{dead} \gets \mathbf{true}$ \Comment{$\Zone$ makes $v$ redundant}
  \EndIf
\EndFor
\While{$\exists$ live $v \in A$ with $\Zone \cup \Zone_v$ convex} \label{ln:union}
  \State $\Zone \gets \Zone \cup \Zone_v$; \quad $v.\mathit{dead} \gets \mathbf{true}$
\EndWhile
\State prepend $\langle \textsc{Intern}(\Zone),\, k,\, \mathit{head},\, \mathbf{false}\rangle$ to $A$ \label{ln:push}
\State \Return the new node
\end{algorithmic}
\end{algorithm}

Line~\ref{ln:bounds} computes the $\aLU{}$ bounds, which depend only on the configuration and
not on the zone. The loop at line~\ref{ln:sub1} decides whether the newcomer adds anything: if
any stored zone already subsumes it the work is finished and we return without touching the
structure, which is the common case in a mature search. Line~\ref{ln:sub2} performs the
converse test, flagging stored zones the newcomer has made redundant --- note that both loops
must run in full, since subsumption is not symmetric and several stored zones may fall to a
single new one. The loop at line~\ref{ln:union} then absorbs any stored zone whose union with
the newcomer is exactly representable, repeating because absorbing one zone can make a further
union convex; it mutates only the incoming zone, never an interned one, which is what makes it
safe while other threads read the table. Finally line~\ref{ln:push} interns the zone and
prepends the node, stamping it with the level so that survivors can be identified later by
integer comparison. Prepending rather than appending is deliberate: it puts the most recently
added zone first, which \cref{sec:batching} uses to detect first touches for free.

The cost of \cref{alg:insert} is dominated by three traversals of $A$, each performing an
$O(m^2)$ matrix comparison per element. The critical section is therefore $\Theta(|A|)$, and
$|A|$ averages $\rho$. When $\rho \approx 1$ this is a formality. When $\rho$ is large it is
the entire program, and the next section is about what to do then.
